\documentclass[11pt]{article}
\usepackage[margin=1in]{geometry}
\usepackage[T1]{fontenc}
\usepackage{lmodern}
\usepackage{amsmath,amssymb}
\usepackage{booktabs,longtable,array}
\usepackage{fancyvrb}
\usepackage[hidelinks]{hyperref}
\setlength{\parindent}{0pt}
\setlength{\parskip}{0.55\baselineskip}
\setlength{\emergencystretch}{3em}
\setcounter{tocdepth}{1}
\setlength{\LTpre}{0pt}
\setlength{\LTpost}{0pt}
\raggedbottom
\makeatletter
\newcommand{\tableinput}[1]{\@@input #1 }
\makeatother
\newcommand{\off}[1]{\texttt{0x#1}}
\title{Static Reconstruction of an Apple Bootloader Image\\
\large \texttt{iboot\_dec\_24a435.bin}}
\author{Binary analysis report}
\date{September 13, 2026}
\begin{document}
\maketitle

\begin{abstract}
The supplied 3,832,328-byte image identifies itself as
\texttt{iBoot for v59}, \texttt{mBoot-20457.2.37}, and \texttt{RELEASE}.
This report reconstructs its startup mechanisms, selected boot-control
paths, command metadata, and subsystem organization from local bytes and
AArch64 disassembly. Supporting artifacts contain a reproducible analyzer,
a linear disassembly of the apparent primary code region, branch candidates,
string references, and command-handler candidates.
\textbf{This is a partial static reconstruction, not recovered original source,
a complete decompilation, or an executable replacement.} No device execution,
firmware modification, or security-policy bypass was performed.
\end{abstract}
\tableofcontents
\newpage

\section{Scope, provenance, and evidence standards}
The only executable input examined is \texttt{iboot\_dec\_24a435.bin}.
Its SHA-256 digest is:
\begin{center}\small\ttfamily
 e1b9aa20c3d00b7eeadecfb1dd04963de9a80\\
 e75830cbc816ede8aed000a0301
\end{center}
The filename's \texttt{24a435} component is not independently authenticated
as an operating-system build number. Likewise, the word \texttt{dec} in the
filename is not proof of its extraction or decryption history. The file
begins with instructions rather than an ELF or Mach-O header. Its readable
code and data are consistent with a raw unpacked bootloader payload.

All hexadecimal locations in this report are \emph{file offsets}, except
where explicitly called physical or linked addresses. Evidence is classified as:
\begin{description}
\item[Observed:] literal bytes, decoded instructions, or an exact branch target.
\item[Corroborated inference:] a role supported by code, references, and repeated metadata patterns.
\item[Unresolved:] a plausible role without a completed control-flow or data-flow proof.
\end{description}
Names such as ``permission gate,'' ``boot coordinator,'' and
\texttt{sub\_00081434} are analytical labels, not original symbols.
Strings alone do not establish that a feature is enabled or reachable.
Apple's general description of iBoot's secure-boot responsibilities provides
context, not proof of this particular build's behavior~\cite{apple}.

\section{Identity and memory layout}
\begin{longtable}{@{}>{\raggedright\arraybackslash}p{0.18\textwidth}>{\raggedright\arraybackslash}p{0.30\textwidth}>{\raggedright\arraybackslash}p{0.45\textwidth}@{}}
\toprule
Location & Observation & Interpretation \\
\midrule
\endhead
\off{00000000} & AArch64 instructions & Entry-like machine setup; little-endian instruction words. \\
\off{00000280} & \texttt{iBoot for v59} & Banner also contains Apple copyright years 2007--2026. \\
\off{000002c0} & \texttt{RELEASE} & Embedded build-style string. \\
\off{00000300} & \texttt{mBoot-20457.2.37} & Embedded build identifier, also referenced by banner code. \\
\off{00000380} & \texttt{0x1fc080000} & Physical relocation destination candidate, loaded at entry. \\
\off{00000408} & \texttt{0xfffffc0000000000} & Address-bias literal used in startup address arithmetic. \\
\off{00219cd8} & Structured pointer data & Apparent end of the main contiguous code region. \\
\off{002b5200} vicinity & Diagnostic strings & Large readable pool, followed by further structured and opaque data. \\
\off{00358000} onward & Data and strings & Secondary metadata/string area; this is not established as executable text. \\
\off{003a7a08} & End of supplied file & Exclusive byte endpoint. \\
\bottomrule
\end{longtable}

The linked-address candidate is
\texttt{0xfffffc01fc080000}. Subtracting it from many 64-bit pointers gives
valid offsets in this image, including pointers to the command names and
startup routines. For such pointers,
\[
 \text{file offset}=\text{linked address}-\text{linked base}.
\]
This rule is not applied blindly to hardware addresses or encoded function
references. In particular, metadata words beginning with \texttt{80f0} or
\texttt{8620} are retained as encoded values rather than dereferenced as
ordinary virtual addresses.

\subsection{Zeroed regions and the file-length discrepancy}
The literal pair at \off{390}/\off{398} corresponds to image-relative
\off{355908}--\off{358000}. Startup clears this region, and its stored
bytes are zero. Another pair at \off{3b8}/\off{3c0} identifies
\off{3a7a40}--\off{3cddd0}, beyond the stored file; this is consistent with
zero-initialized runtime storage.
The rounded relocation extent \off{3a7a40} is 56 bytes beyond the file's
\off{3a7a08} endpoint. Loader-supplied padding or alignment could explain
this, but the required initial memory contents are not recoverable from the
file alone. The image must not be treated as a complete runtime memory dump.

\section{Startup reconstruction}
\subsection{Entry, relocation, and overlap handling}
The initial sequence writes \texttt{VTTBR\_EL2} and updates
\texttt{HCR\_EL2}, with synchronization barriers between stages.
At \off{1c}, \texttt{ADRP}/\texttt{ADD} derives the current image location;
\off{24} loads the preferred destination. A direct call at \off{28}
targets \off{3b338}, a platform helper whose complete semantics remain unresolved.
The comparison at \off{2c} branches to \off{18c} when relocation is unnecessary.

When relocation is needed, the code computes a 64-byte-aligned copy length
from header literals. Address comparisons at \off{68}--\off{84} select
safe handling for overlapping ranges. A small copy routine is staged in
scratch memory before execution transfers to it:
\begin{itemize}
\item \off{f0}--\off{124}: forward copy, 64 bytes per iteration.
\item \off{128}--\off{15c}: reverse copy using pre-decremented end pointers.
\item \off{160}--\off{174}: erase a source-region remainder in 64-byte chunks.
\item \off{178}--\off{184}: erase the staged scratch code in 16-byte chunks.
\item \off{188}: indirect transfer to the destination image.
\end{itemize}
Cache invalidation and synchronization appear around the copy routines.
This is more specific than a generic ``load firmware'' operation: the entry
actively relocates itself and removes selected temporary copies.

\subsection{Runtime establishment}
At \off{18c}, the code masks interrupts, writes an implementation-specific
system register, and sets \texttt{VBAR\_EL1} from an address corresponding
to offset \off{1b6800}. It clears three literal-delimited memory ranges,
then branches from \off{1f8} to \off{33180}.
That continuation resets \texttt{TTBR0\_EL1}/\texttt{TTBR1\_EL1}, establishes
a stack, calls several platform routines, and performs further processor
configuration. The helpers implement additional mapping and initialization;
the full hardware register model has not been recovered.

The sequence at \off{33248}--\off{33280} is an initializer-dispatch loop:
it reads 32-bit entries, adds a base, applies \texttt{PACIZA}, and calls
through \texttt{BLRAAZ}. Crucially, the start and end literals at
\off{33288}/\off{33290} are \emph{equal} in this image, so that loop executes
zero callbacks for these stored values. Its existence does not imply that
initializers actually run. The next direct transfer is
\off{33284} $\rightarrow$ \off{1bdb98}; the intervening runtime-to-boot
initialization chain has not been fully reconstructed.

\section{Boot coordination and recovery}
The routine near \off{4080} is a strong boot-coordinator candidate.
It references the banner, board/revision format, build tag, build style,
and serial-number format. It also references recovery settings and directly
calls the autoboot routine at \off{486c} $\rightarrow$ \off{4e6c}.
The following anchors establish actual code references rather than merely
words found in the binary:
\begin{longtable}{@{}llp{0.45\textwidth}@{}}
\toprule
Code offset & String offset & Referenced text \\
\midrule\endhead
\off{47fc} & \off{2b55d0} & \texttt{boot-command} \\
\off{4804} & \off{2b5386} & \texttt{device-recovery} \\
\off{4890} & \off{2b5445} & \texttt{delay-recovery-image} \\
\off{4974} & \off{2b546a} & Entering recovery mode, starting command prompt \\
\off{4d18} & \off{2b5661} & \texttt{recover-once} \\
\off{4d48} & \off{2b5646} & \texttt{one-time-boot-command} \\
\off{4ed4} & \off{2b55fe} & \texttt{auto-boot} \\
\off{4f94} & \off{2b5610} & \texttt{bootdelay} \\
\off{5018} & \off{2b561a} & aborting autoboot due to user intervention \\
\bottomrule
\end{longtable}

A conservative behavioral model is: establish the runtime and device state;
consult boot/recovery settings; attempt the configured boot path when allowed;
otherwise enter the recovery-command environment. This is a high-level
inference, not a proof of the ordering of every policy check. In particular,
one-shot variables, upgrade handling, and failure fallback precedence require
further tracing.

\subsection{Autoboot countdown: a verified local behavior}
Within the path that reaches \off{4fb8}, the returned \texttt{bootdelay}
value is retained and multiplied by 1,000,000 using 32-bit arithmetic.
A polling loop calls \off{7fd28}, tests the result for 13 or 10
(carriage return or line feed), compares an elapsed counter against the
threshold, and calls \off{1b65e0} with 10,000 between polls.
The counter and delay units are not established here, although the constants
are consistent with microsecond-based timing. On either tested key value,
the loop prints the intervention message and leaves the countdown path.
The timeout comparison is unsigned and strict \texttt{greater-than}.

\section{Recovered command metadata}
The image contains repeated records with duplicate command-name pointers,
a common descriptor pointer, and an encoded handler word. For example,
\off{21b388} references \texttt{bootx}; its encoded word is
\texttt{0x80f000000000121c}, whose low 32 bits identify plausible handler
code at \off{121c}. Adjacent records commonly have a stride of 120 bytes.
The \texttt{go} record uses a different upper encoding,
\texttt{0x86200000000023b0}; its low-bit target is corroborated by memory-image
validation and permission-failure code.

The table is a \emph{candidate inventory}, not a proof that all names are
exposed in every boot mode, nor a complete decode of the record ABI.
\begin{longtable}{@{}lll@{}}
\toprule
Name & Record offset & Handler offset \\
\midrule\endhead
\tableinput{reconstruction/commands.tex}
\bottomrule
\end{longtable}

\subsection{Handlers with code-level corroboration}
\textbf{\texttt{memboot}:} the handler reads \texttt{filesize} using a helper
at \off{16190c}. Zero produces the invalid-or-missing diagnostic.
An unsigned comparison rejects sizes at or above \off{20000001}, so the
accepted nonzero size range at this local check is 1--536,870,912 bytes
(512 MiB maximum). It then calls \off{81434} and tests result bit 0;
a clear bit leads to ``Permission Denied.'' Subsequent comparisons verify
that the requested span fits within returned bounds before further processing.
These are observed local checks, not a complete verification policy.

\textbf{\texttt{bootx}:} the handler calls the same gate and has a matching
permission-denied branch. It calls the BootKC-related routine at
\off{29220}, follows multiple negative-result branches, and updates
\texttt{boot-breadcrumbs}. The final transfer mechanism and all image
validation steps are not fully reconstructed.

\textbf{\texttt{go}:} the handler calls \off{81434}, then a separate routine
at \off{c6508}. A nonzero result from the latter reaches ``Memory image not
valid''; permission failure is a distinct branch. Thus the name \texttt{go}
is not evidence of an unrestricted jump primitive.

\textbf{\texttt{firmware}:} the default-argument path checks \texttt{filesize}
and calls the shared gate before continuing. Its alternate-argument path
uses other helpers and has not been fully characterized.

\textbf{\texttt{reboot}/\texttt{reset}:} both records point to
\off{5628}, establishing an alias at the candidate-handler level.
\textbf{\texttt{saveenv}:} \off{5624} is a one-instruction branch to
\off{1623f4}, not a standalone persistence implementation.

\textbf{Environment commands:} \texttt{getenv}, \texttt{setenv},
\texttt{setenvnp}, and \texttt{clearenvp} lead into the \off{162000}
region. Names suggest environment management, but exact persistent versus
nonpersistent semantics are not established merely by their suffixes.

\textbf{Display/ramdisk commands:} \texttt{bgcolor}, \texttt{setpicture},
\texttt{devicetree}, and \texttt{ramdisk} have distinct candidate handlers.
The \texttt{setpicture} usage strings describe memory, image-type, and
file-based inputs and masks for update, preservation, restore mode, and
flipbook selection. The image additionally contains corruption and
permission diagnostics referenced by that handler.

\section{Subsystem reconstruction}
\begin{longtable}{@{}>{\raggedright\arraybackslash}p{0.20\textwidth}>{\raggedright\arraybackslash}p{0.34\textwidth}>{\raggedright\arraybackslash}p{0.37\textwidth}@{}}
\toprule
Subsystem & Local evidence & Supported conclusion and limit \\
\midrule\endhead
Boot objects & \texttt{BootKC-*} references at \off{2a808}--\off{2a900}; \texttt{BootKC} at \off{293fc} & Separate labels for read-only, executable, writable, virtual-address, and entry metadata. Exact object layout remains unresolved. \\
Manifests & \off{6ae0} references ``verify boot manifest hash''; \off{1f9c90} references a manifest-decode error & Manifest handling exists. The latter is not independently identified as the secure-boot verifier; subsystem ownership matters. \\
Secure payloads & Names for SecurePageTableMonitor, TrustedExecutionMonitor, cL4, SEP firmware, and trust caches near \off{2b5ac4} & Loadable-payload vocabulary is present. This does not prove all payloads are used on v59. \\
SPTM layout & \off{2f5d0} and \off{304c4} reference \texttt{lay\_out\_sptm\_objects} & Code-level anchors for a secure-page-table-monitor layout path. Full semantics are not recovered. \\
Storage & \off{76fe8} references \texttt{nvme\_nand\%d}; additional firmware/syscfg device names & NVMe-oriented storage naming and initialization are represented. A full filesystem or controller implementation has not been reconstructed. \\
USB recovery & \off{c8d1c} references ``usb rpc task''; recovery-device and USB-serial strings & USB task/recovery infrastructure is present. The protocol, packet formats, and exposure conditions remain unresolved. \\
Memory training & \off{53ae4} references memory CA calibration failure; extensive Vref and channel diagnostics & Platform memory-calibration code is represented; no full DRAM training algorithm is claimed. \\
Cryptographic helpers & \off{88e04}/\off{88e14} reference AES argument/key-size errors & AES-related argument validation is present. Algorithms, modes, keys, and secure-boot binding have not been fully traced. \\
Runtime checks & \off{1f54f4} references \texttt{\_\_firebloom\_panic}; pointer bounds comparisons and authenticated returns & Evidence of runtime checks and pointer authentication in sampled code, not a proof of memory safety or exploitability. \\
Power and diagnostics & Battery state, voltage, panic log, and poweroff strings with platform metadata & Relevant support paths are present; hardware effects require further reverse engineering. \\
\bottomrule
\end{longtable}

\section{Reconstructed pseudocode}
The accompanying \texttt{reconstruction/behavior.pseudo} expresses only
three bounded observations: relocation stages, the \texttt{memboot} validation
prefix, and the countdown loop. Calls retain offset-based names when their
semantics are not proven. It intentionally stops before unresolved processing
rather than inventing parsers, cryptographic checks, or a kernel handoff.
\VerbatimInput[fontsize=\footnotesize]{reconstruction/behavior.pseudo}

\section{Reproducibility and coverage}
Run the analyzer from the workspace directory:
\begin{Verbatim}[fontsize=\small]
python reconstruct.py
latexmk -pdf -interaction=nonstopmode -halt-on-error main.tex
\end{Verbatim}
The analyzer uses Python's standard library and the installed LLVM 19 shared
library through its C disassembler interface~\cite{llvm}. No additional
packages or virtual environment are needed. Its input hash is checked before
analysis. Address bases, code boundary, and command encodings are explicitly
image-specific, not a universal iBoot parser.

\begin{center}
\begin{tabular}{@{}lr@{}}
\toprule
Inventory metric & Count \\
\midrule
\tableinput{reconstruction/metrics.tex}
\bottomrule
\end{tabular}
\end{center}
These numbers are \emph{scan results}, not numbers of verified functions or
reachable instructions. Printable runs include accidental matches inside
code and opaque data. Candidate starts combine direct-call targets and
pointer-authentication prologues, including thunks and internal entries.
The 692 undecoded words include literal pools or unsupported encodings;
not all represent missing instructions. Conversely, successful decoding
alone does not prove that a word is executable.

\subsection{Artifact guide}
\begin{description}
\item[\texttt{linear\_disassembly.txt.gz}:] AArch64 linear sweep through
\off{219cd8}, with candidate labels and selected annotations. Embedded
header/data bytes are included and must not be mistaken for code.
\item[\texttt{selected\_disassembly.txt}:] Readable slices for startup,
boot coordination, countdown, and candidate command handlers.
\item[\texttt{commands.json}:] Raw encoded handler values, name locations,
metadata record locations, and inferred target offsets.
\item[\texttt{branches.json}:] Direct branch/call candidates; indirect
calls are not resolved. ``Nearest candidate'' is proximity, not function ownership.
\item[\texttt{string\_xrefs.json}, \texttt{address\_xrefs.json}:]
Short-window \texttt{ADRP}/\texttt{ADD} matches, not full data-flow analysis.
\item[\texttt{pointers.json}, \texttt{strings.tsv}:] Aligned linked-pointer
matches and printable byte runs across the whole image.
\item[\texttt{candidate\_starts.json}, \texttt{entropy.json}:]
Entry hypotheses and 64-KiB block statistics. Entropy is not a compression
or encryption classifier.
\end{description}
All listed artifacts are under \texttt{reconstruction/}. The analyzer is
\texttt{reconstruct.py}; the report source is \texttt{main.tex}.

\section{What is still required for a full reconstruction}
A complete behavioral reconstruction would require verified function
boundaries, indirect-call resolution, structure/type recovery, and
path-sensitive analysis of the boot-policy and image-authentication routines.
The opaque middle payloads and possible embedded firmware require separate
format identification. The present sweep does not disassemble them as
independent processors' firmware.

The key next targets are the shared gate \off{81434}, image-processing
routine \off{c6508}, the environment/command dispatcher, the runtime bridge
from \off{1bdb98}, and the actual kernel or monitor handoff. Calling
conventions must be recovered before replacing the multi-register pointer
representations with ordinary C pointers. Validation against an appropriate
hardware or emulation environment would then be necessary; no such validation
was performed here.

\textbf{Bottom line:} the supplied bytes support a substantial indexed
static-analysis corpus and several concrete local reconstructions. They do
not yet support claiming a full source recovery, a verified clone, or a
complete account of this binary's behavior.

\begin{thebibliography}{9}
\bibitem{apple} Apple, \emph{Apple Platform Security: Boot process for a Mac
with Apple silicon}, published February 18, 2021; consulted September 13,
2026. General platform context only, not identification of the supplied board.
\bibitem{llvm} LLVM Project, \emph{Disassembler.h: C Disassembler Interface},
\texttt{LLVMCreateDisasmCPUFeatures}, \texttt{LLVMDisasmInstruction}, and
\texttt{LLVMDisasmDispose}; consulted September 13, 2026. Local implementation:
installed LLVM 19 library.
\end{thebibliography}
\end{document}
